% CSE 715 — Computer Graphics
% Chapter 11: Color & Shading Models (Phong, RGB interpolation, monitor matrix M, point left/right test, color gamut)
% University of Chittagong, 7th Semester B.Sc. Engineering
% Source: Schaum's Outline of Computer Graphics (Xiang & Plastock), Chapter 11
%
% Compile with: pdflatex Chapter11_Solutions.tex (run twice for TOC)

\documentclass[12pt, a4paper]{article}

\usepackage[left=2.5cm, right=2.5cm, top=2.8cm, bottom=2.8cm, headheight=15pt]{geometry}
\usepackage{amsmath, amssymb, mathtools}
\usepackage{tikz}
\usetikzlibrary{arrows.meta, positioning, shapes.geometric, calc, backgrounds, fit}
\usepackage{booktabs}
\usepackage{array}
\usepackage{xcolor}
\usepackage{enumitem}
\usepackage{fancyhdr}
\usepackage{titlesec}
\usepackage{mdframed}
\usepackage{tabularx}
\usepackage{hyperref}

\definecolor{sectionblue}{RGB}{10,60,110}
\definecolor{boxbg}{RGB}{242,247,255}
\definecolor{answerbg}{RGB}{240,255,240}
\definecolor{notesbg}{RGB}{255,250,230}

\mdfdefinestyle{solutionframe}{linecolor=sectionblue, backgroundcolor=boxbg, roundcorner=4pt, linewidth=1.4pt, innertopmargin=7pt, innerbottommargin=7pt, innerleftmargin=9pt, innerrightmargin=9pt}
\mdfdefinestyle{answerframe}{linecolor=green!55!black, backgroundcolor=answerbg, roundcorner=4pt, linewidth=1.4pt, innertopmargin=7pt, innerbottommargin=7pt, innerleftmargin=9pt, innerrightmargin=9pt}
\mdfdefinestyle{notesframe}{linecolor=orange!70!black, backgroundcolor=notesbg, roundcorner=4pt, linewidth=1pt, innertopmargin=5pt, innerbottommargin=5pt, innerleftmargin=8pt, innerrightmargin=8pt}
\newenvironment{solutionbox}{\begin{mdframed}[style=solutionframe]}{\end{mdframed}}
\newenvironment{answerbox}{\begin{mdframed}[style=answerframe]}{\end{mdframed}}
\newenvironment{notesbox}{\begin{mdframed}[style=notesframe]}{\end{mdframed}}

\pagestyle{fancy}
\fancyhf{}
\lhead{\small\color{sectionblue} CSE 715 --- Chapter 11: Color \& Shading Models}
\rhead{\small\color{sectionblue} Model Solutions (2020--2024)}
\cfoot{\small\thepage}
\renewcommand{\headrulewidth}{0.5pt}

\titleformat{\section}{\Large\bfseries\color{sectionblue}}{}{0em}{}[\vspace{-2pt}\color{sectionblue}\rule{\linewidth}{0.8pt}]
\titleformat{\subsection}{\large\bfseries\color{sectionblue!85!black}}{}{0em}{}

\begin{document}

\begin{center}
  {\LARGE\bfseries\color{sectionblue} CSE 715: Computer Graphics}\\[0.6em]
  {\Large\bfseries Chapter 11 --- Color \& Shading Models}\\[0.3em]
  {\large Phong Model, RGB Interpolation, Monitor Color Matrix $M$, Point Left/Right Test, Color Gamut}\\[0.2em]
  {\normalsize University of Chittagong \quad|\quad 7\textsuperscript{th} Semester B.Sc.\ (Engg.) Examination}\\[0.3em]
  {\small Source: Schaum's Outline of Computer Graphics (Xiang \& Plastock), Chapter 11 --- this is the ``Q7 cluster'' (Section B) every year, occasionally spilling into Q6/Q8}
\end{center}

\vspace{0.3em}\noindent\rule{\linewidth}{1pt}\vspace{0.5em}

\begin{notesbox}
\textbf{Yield note (re-verified against all 5 raw scanned papers 2026-07-07):} this chapter bundles \textbf{5 subtopics} that all live in the same exam-question cluster: RGB scanline interpolation (2021 2022 2023), Monitor color matrix $M$ (2021 2023), Phong illumination model (2022 2023), Point left/right-of-line test (2020 2022), Color gamut (2022). \textbf{Every one of the 5 years has at least 2--3 of these}, so this is a real, guaranteed-to-recur Section-B question, not a stretch-goal. 2024 does not touch this cluster at all (its Q7/Q8 go to visible-surface/trapezoid-fill/ray-tracing instead). \textbf{Correction:} the Phong-model year attribution in \texttt{\_TopicQuestionMap.md} previously said "2022 Q6d" --- re-reading the raw 2022 scan confirms it is actually \textbf{Q7e}, not Q6d. Fixed in this pass.
\end{notesbox}

\tableofcontents
\newpage

%=====================================================================
\section{Theoretical Framework}
%=====================================================================

\subsection{The Phong Illumination Model (\S11.2)}
\begin{answerbox}
For a point $Q$ on a surface, light source direction $\mathbf{L}$, surface normal $\mathbf{N}$, mirror-reflection direction $\mathbf{R}$, and viewer direction $\mathbf{V}$ (all unit vectors), the total reflected intensity is
$$I = I_a k_a + I_p\big(k_d\cos\theta + k_s\cos^{n}\varphi\big) = I_a k_a + I_p\big(k_d\,\mathbf{L}\!\cdot\!\mathbf{N} + k_s(\mathbf{R}\!\cdot\!\mathbf{V})^{n}\big)$$
where $\theta$ is the angle between $\mathbf{L}$ and $\mathbf{N}$ (angle of incidence), $\varphi$ is the angle between $\mathbf{R}$ and $\mathbf{V}$, and $0\le k_a,k_d,k_s\le1$ are the surface's ambient/diffuse/specular reflection coefficients. $\mathbf{R} = 2(\mathbf{L}\!\cdot\!\mathbf{N})\mathbf{N}-\mathbf{L}$.
\begin{itemize}[noitemsep]
  \item \textbf{Ambient term} $I_ak_a$: directionless — a flat "fill" so unlit surfaces aren't pure black (accounts for indirect/bounced light).
  \item \textbf{Diffuse term} $I_p k_d\cos\theta$: models a \emph{matte} surface scattering light equally in all directions; brightness depends only on the incidence angle $\theta$ (Lambert's law), \textbf{not} on viewer position — this is what gives an object its basic 3D "shape shading."
  \item \textbf{Specular term} $I_p k_s\cos^{n}\varphi$: models the \emph{shiny highlight} on glossy surfaces, concentrated tightly around the mirror direction $\mathbf{R}$; \textbf{does} depend on viewer position $\mathbf{V}$. The shininess exponent $n$ controls the highlight's tightness ($n=1$: dull, $n=100$: near-mirror).
\end{itemize}
For multiple light sources, the diffuse+specular term sums over all sources $i$: $I = I_ak_a + \sum_i I_{p_i}(k_d\mathbf{L}_i\!\cdot\!\mathbf{N} + k_s(\mathbf{R}_i\!\cdot\!\mathbf{V})^n)$.
\end{answerbox}

\subsection{Interpolative Shading Methods (\S11.3)}
\begin{answerbox}
Evaluating the Phong formula at every pixel is expensive, so real renderers evaluate it only at selected points and interpolate elsewhere:
\begin{itemize}[noitemsep]
  \item \textbf{Constant shading:} evaluate once per polygon (e.g.\ at its centroid), flat-fill the whole face. Fast, but false contours appear at polygon edges and specular highlights are often lost/distorted.
  \item \textbf{Gouraud shading:} evaluate the Phong formula at each \emph{vertex}, then bilinearly interpolate the resulting \emph{colors} across the polygon (via the two-edge-intersection method: interpolate along each edge to the scanline, then interpolate between those two points). Removes false contours even on a coarse mesh, but a fine mesh is still needed to capture sharp specular highlights.
  \item \textbf{Phong shading:} interpolate the \emph{normal vectors} $\mathbf{N}$ (not the colors) across the polygon, then evaluate the full Phong formula at every pixel using the interpolated normal. More expensive, but handles specular highlights far better.
\end{itemize}
Vertex normals for a polygon mesh are obtained either from the underlying curved surface (true geometric normal) or, absent one, by \textbf{averaging the normals of all polygons meeting at that vertex} (then renormalizing).
\end{answerbox}

\subsection{RGB Linear Interpolation Along a Scanline}
\begin{answerbox}
Given a point $P_1$ with color/intensity $I_1$ on line $y_1$ and $P_2$ with color $I_2$ on line $y_2$, the color at any line $y'$ is found by linear interpolation, applied \textbf{component-wise} to $(R,G,B)$:
$$I' = I_1 + (I_2-I_1)\cdot\frac{y'-y_1}{y_2-y_1}$$
This is the same formula whether $y'$ falls \emph{between} $y_1,y_2$ (interpolation) or outside that range (extrapolation) — the fraction $\frac{y'-y_1}{y_2-y_1}$ just falls outside $[0,1]$ in the extrapolation case.
\end{answerbox}

\subsection{Monitor Color Transformation Matrix $M$ (CIE $XYZ$)}
\begin{answerbox}
\textbf{Goal:} given the monitor's white point $(x_w,y_w,Y_w)$ (the $XYZ$ chromaticity produced when $R=G=B=1$) and the chromaticity coordinates $(x_r,y_r),(x_g,y_g),(x_b,y_b)$ of its three phosphors, find $M$ such that $\begin{psmallmatrix}X\\Y\\Z\end{psmallmatrix}=M\begin{psmallmatrix}R\\G\\B\end{psmallmatrix}$.
\textbf{Step 1 --- white point to $XYZ$:} $X_w = \dfrac{x_w}{y_w}Y_w,\quad Z_w=\dfrac{1-x_w-y_w}{y_w}Y_w$.
\textbf{Step 2 --- auxiliary $z$'s:} $z_r=1-x_r-y_r,\ z_g=1-x_g-y_g,\ z_b=1-x_b-y_b$.
\textbf{Step 3 --- solve for $C_r,C_g,C_b$} from
$$\begin{pmatrix}X_w\\Y_w\\Z_w\end{pmatrix}=\begin{pmatrix}x_r&x_g&x_b\\y_r&y_g&y_b\\z_r&z_g&z_b\end{pmatrix}\begin{pmatrix}C_r\\C_g\\C_b\end{pmatrix}$$
\textbf{Step 4 --- assemble $M$:}
$$M=\begin{pmatrix}x_rC_r&x_gC_g&x_bC_b\\y_rC_r&y_gC_g&y_bC_b\\z_rC_r&z_gC_g&z_bC_b\end{pmatrix}$$
\textbf{Sanity check:} each row of $M$ must sum to $X_w,Y_w,Z_w$ respectively (since $M\cdot(1,1,1)^T=(X_w,Y_w,Z_w)^T$ when $R=G=B=1$).
\end{answerbox}

\subsection{Color Gamut of a Monitor}
\begin{answerbox}
The \textbf{color gamut} of a monitor is the range of colors it can physically display. On the CIE chromaticity diagram, the chromaticity coordinates $(x_r,y_r),(x_g,y_g),(x_b,y_b)$ of the R/G/B phosphors form the three vertices of a \textbf{triangle}; any color \emph{inside} that triangle can be produced as some non-negative combination of $R,G,B$, while more saturated colors \emph{outside} the triangle cannot be reproduced exactly by that monitor. Different monitors (different phosphors) have different, generally non-identical gamut triangles.
\end{answerbox}

\subsection{Point Left/Right-of-Line Test}
\begin{answerbox}
Given a directed line segment from $P_1(x_1,y_1)$ to $P_2(x_2,y_2)$ and a test point $P(x,y)$, define
$$C = (x_2-x_1)(y-y_1) - (y_2-y_1)(x-x_1)$$
This is the $z$-component of the 2D cross product $\overrightarrow{P_1P_2}\times\overrightarrow{P_1P}$. In a standard right-handed $xy$ system:
\begin{itemize}[noitemsep]
  \item $C>0$: $P$ lies to the \textbf{left} of the directed line $P_1\to P_2$.
  \item $C<0$: $P$ lies to the \textbf{right} of the directed line $P_1\to P_2$.
  \item $C=0$: $P$ lies \textbf{on} the line.
\end{itemize}
This single sign test is the building block for polygon-orientation checks, convex-polygon clipping (Sutherland-Hodgman "inside" tests), and back-face determination.
\end{answerbox}

%=====================================================================
\section{Examination 2020 --- Q7(a): Point Left/Right Test [2.25 marks]}
%=====================================================================
\begin{solutionbox}
\textbf{Question:} Show the conditions to identify whether a point locates the right or left side of a line segment.
\end{solutionbox}
\begin{answerbox}
Derive $C=(x_2-x_1)(y-y_1)-(y_2-y_1)(x-x_1)$ as in \S1.5 (cross-product $z$-component of $\overrightarrow{P_1P_2}\times\overrightarrow{P_1P}$), then state: $C>0\Rightarrow$ left, $C<0\Rightarrow$ right, $C=0\Rightarrow$ on the line.
\end{answerbox}

%=====================================================================
\section{Examination 2021 --- Q7(a) \& Q7(c)}
%=====================================================================

\subsection*{Q7(a) --- Monitor Color Matrix $M$ [3 marks]}
\begin{solutionbox}
\textbf{Question:} Assume that a monitor produces the standard white $D_{65}$ with $x_w=0.313$, $y_w=0.329$, $Y_w=1.0$ when $R=G=B=1$, and the chromaticity coordinates of its phosphors are $R(0.62,0.34)$, $G(0.29,0.59)$, $B(0.15,0.06)$. Find the color transformation matrix $M$ for the monitor.
\end{solutionbox}
\begin{answerbox}
$X_w = 0.313/0.329 = 0.951$, $Z_w=(1-0.313-0.329)/0.329=1.088$. Auxiliary $z$'s: $z_r=0.04,\ z_g=0.12,\ z_b=0.79$. Solving
$$\begin{pmatrix}0.951\\1.0\\1.088\end{pmatrix}=\begin{pmatrix}0.62&0.29&0.15\\0.34&0.59&0.06\\0.04&0.12&0.79\end{pmatrix}\begin{pmatrix}C_r\\C_g\\C_b\end{pmatrix}\ \Rightarrow\ C_r=0.705,\ C_g=1.170,\ C_b=1.164$$
$$M=\begin{pmatrix}x_rC_r&x_gC_g&x_bC_b\\y_rC_r&y_gC_g&y_bC_b\\z_rC_r&z_gC_g&z_bC_b\end{pmatrix}=\begin{pmatrix}0.437&0.339&0.175\\0.240&0.690&0.070\\0.028&0.140&0.920\end{pmatrix}$$
\textbf{Check:} row sums $=0.951,\ 1.0,\ 1.088$ \checkmark matches $X_w,Y_w,Z_w$.
\end{answerbox}

\subsection*{Q7(c) --- RGB Scanline Interpolation [2 marks]}
\begin{solutionbox}
\textbf{Question:} If point $P_1$ on line 6 has an RGB color $(1,1,0)$ and point $P_2$ on line 15 has an RGB color $(0.6,0.5,0.2)$, what is the color for line 8?
\end{solutionbox}
\begin{answerbox}
Fraction $t=\dfrac{8-6}{15-6}=\dfrac{2}{9}=0.222$.
$$R = 1+(0.6-1)(0.222) = 0.911 \qquad G = 1+(0.5-1)(0.222)=0.889 \qquad B = 0+(0.2-0)(0.222)=0.044$$
\textbf{Color at line 8} $\approx (0.911,\ 0.889,\ 0.044)$.
\end{answerbox}

%=====================================================================
\section{Examination 2022 --- Q7(a), Q7(d), Q7(e) \& Q8(b)}
%=====================================================================

\subsection*{Q7(a) --- Point Left/Right Test [1 mark]}
\begin{solutionbox}
\textbf{Question:} $C=(x_2-x_1)(y-y_1)-(y_2-y_1)(x-x_1)$. Refer to this expression, what are the conditions whether a point is left or right of a line segment?
\end{solutionbox}
\begin{answerbox}
$C>0 \Rightarrow$ point is left of the directed segment $P_1\to P_2$; $C<0\Rightarrow$ right; $C=0\Rightarrow$ on the line (see \S1.5).
\end{answerbox}

\subsection*{Q7(d) --- RGB Scanline Interpolation [2.5 marks]}
\begin{solutionbox}
\textbf{Question:} If point $p_1$ on line 13 has an RGB color $(0.4,5,0)$ and point $p_2$ on line 15 has an RGB color $(2,0.5,7)$, what is the color for line 12?
\end{solutionbox}
\begin{answerbox}
Line 12 lies \emph{outside} $[13,15]$, so this is a linear \textbf{extrapolation} — same formula, fraction $t=\dfrac{12-13}{15-13}=-0.5$.
$$R=0.4+(2-0.4)(-0.5)=-0.4 \qquad G=5+(0.5-5)(-0.5)=7.25 \qquad B=0+(7-0)(-0.5)=-3.5$$
\textbf{Color at line 12} $=(-0.4,\ 7.25,\ -3.5)$. (Values fall outside $[0,1]$ because the given endpoint values themselves already exceed the normalized range — the question tests the interpolation \emph{mechanism}, not physical color validity.)
\end{answerbox}

\subsection*{Q7(e) --- Significance of Diffuse and Specular Reflection [1.5 marks]}
\begin{solutionbox}
\textbf{Question:} Enumerate the significance of diffuse and specular reflection in the Phong illumination model.
\end{solutionbox}
\begin{answerbox}
\begin{itemize}[noitemsep]
  \item \textbf{Diffuse reflection} ($k_d\cos\theta$): models light scattered equally in all directions off a matte surface (Lambertian); intensity depends only on the incidence angle $\theta=\angle(\mathbf{L},\mathbf{N})$, \emph{not} on where the viewer is standing. This is what gives an object its basic, direction-independent 3D "form shading."
  \item \textbf{Specular reflection} ($k_s\cos^n\varphi$): models the shiny highlight seen on glossy surfaces, concentrated near the mirror-reflection direction $\mathbf{R}$; intensity depends on $\varphi=\angle(\mathbf{R},\mathbf{V})$, so it \emph{does} shift with viewer position. The exponent $n$ (shininess) controls how tight/mirror-like the highlight looks.
  \item Together with the ambient term, they let a purely local (non-global) illumination model produce a visually convincing "shape + shiny highlight" appearance cheaply.
\end{itemize}
\end{answerbox}

\subsection*{Q8(b) --- Color Gamut of the Monitor [1 mark]}
\begin{solutionbox}
\textbf{Question:} Explain the color Gamut of the monitor.
\end{solutionbox}
\begin{answerbox}
See \S1.6: the triangle in the CIE chromaticity diagram whose vertices are the monitor's R/G/B phosphor chromaticities; colors inside are displayable, colors outside (more saturated) are not.
\end{answerbox}

%=====================================================================
\section{Examination 2023 --- Q6(d), Q7(b) \& Q7(c)}
%=====================================================================

\subsection*{Q6(d) --- Explain the Phong Diagram [3 marks]}
\begin{solutionbox}
\textbf{Question:} In light of the Phong model, explain the following image (light source, point $Q$, vectors $\mathbf{L}$, $\mathbf{N}$, $\mathbf{R}$, $\mathbf{V}$, angles $\theta,\theta,\varphi$).
\end{solutionbox}
\begin{answerbox}
\begin{itemize}[noitemsep]
  \item $\mathbf{L}$: direction from surface point $Q$ toward the light source.
  \item $\mathbf{N}$: surface normal at $Q$.
  \item $\mathbf{R}$: direction of perfect mirror reflection of $\mathbf{L}$ about $\mathbf{N}$ — by the law of reflection the angle between $\mathbf{L}$ and $\mathbf{N}$ equals the angle between $\mathbf{R}$ and $\mathbf{N}$ (both labeled $\theta$ in the figure).
  \item $\mathbf{V}$: direction from $Q$ toward the viewpoint/eye.
  \item $\varphi$: the angle between $\mathbf{R}$ and $\mathbf{V}$ — the smaller $\varphi$ is (viewer close to the mirror-reflection direction), the stronger the specular highlight, per $\cos^n\varphi$.
\end{itemize}
Together the figure is the geometric setup behind $I=I_ak_a+I_p(k_d\,\mathbf{L}\!\cdot\!\mathbf{N}+k_s(\mathbf{R}\!\cdot\!\mathbf{V})^n)$ from \S1.1: the diffuse term depends only on $\theta$ (viewer-independent), the specular term depends on $\varphi$ (viewer-dependent).
\end{answerbox}

\subsection*{Q7(b) --- Monitor Color Matrix $M$ [3 marks]}
\begin{solutionbox}
\textbf{Question:} Assume that a monitor produces the standard white $D_{65}$ with $x_w=0.3$, $y_w=0.2$, $Y_w=1.0$ when $R=G=B=1$, and the chromaticity coordinates of its phosphors are $R(0.62,0.34)$, $G(0.29,0.59)$, $B(0.15,0.06)$. Find the color transformation matrix $M$.
\end{solutionbox}
\begin{answerbox}
$X_w=0.3/0.2=1.5,\quad Z_w=(1-0.3-0.2)/0.2=2.5$. Same phosphors as \S3 (2021), but different white point, so solve fresh:
$$\begin{pmatrix}1.5\\1.0\\2.5\end{pmatrix}=\begin{pmatrix}0.62&0.29&0.15\\0.34&0.59&0.06\\0.04&0.12&0.79\end{pmatrix}\begin{pmatrix}C_r\\C_g\\C_b\end{pmatrix}\ \Rightarrow\ C_r\approx1.427,\ C_g\approx0.567,\ C_b\approx3.007$$
$$M\approx\begin{pmatrix}0.885&0.164&0.451\\0.485&0.335&0.180\\0.057&0.068&2.375\end{pmatrix}$$
\textbf{Check:} row sums $\approx1.500,\ 1.000,\ 2.500$ \checkmark matches $X_w,Y_w,Z_w$.
\end{answerbox}

\subsection*{Q7(c) --- RGB Scanline Interpolation [2 marks]}
\begin{solutionbox}
\textbf{Question:} If point $p_1$ on line 3 has an RGB color $(5,5,0)$ and point $p_2$ on line 9 has an RGB color $(2,0.5,7)$, what is the color for line 5?
\end{solutionbox}
\begin{answerbox}
Fraction $t=\dfrac{5-3}{9-3}=\dfrac{1}{3}$.
$$R=5+(2-5)\tfrac13=4 \qquad G=5+(0.5-5)\tfrac13=3.5 \qquad B=0+(7-0)\tfrac13=2.333$$
\textbf{Color at line 5} $=(4,\ 3.5,\ 2.333)$.
\end{answerbox}

%=====================================================================
\section{Quick Summary}
%=====================================================================
\begin{center}
\renewcommand{\arraystretch}{1.3}
\begin{tabular}{p{4.6cm}|p{2.4cm}|p{6.9cm}}
\toprule
\textbf{Subtopic} & \textbf{Years} & \textbf{Method} \\
\midrule
RGB scanline interpolation & 2021, 2022, 2023 & $I'=I_1+(I_2-I_1)\frac{y'-y_1}{y_2-y_1}$, component-wise; works for extrapolation too \\
Monitor color matrix $M$ & 2021, 2023 & $X_w,Z_w$ from white point $\to$ solve 3$\times$3 for $C_r,C_g,C_b$ $\to$ scale each phosphor column \\
Phong illumination model & 2022, 2023 & $I=I_ak_a+I_p(k_d\mathbf{L}\!\cdot\!\mathbf{N}+k_s(\mathbf{R}\!\cdot\!\mathbf{V})^n)$ — diffuse viewer-independent, specular viewer-dependent \\
Point left/right test & 2020, 2022 & $C=(x_2-x_1)(y-y_1)-(y_2-y_1)(x-x_1)$; $C{>}0$ left, $C{<}0$ right, $C{=}0$ on line \\
Color gamut & 2022 & Triangle of phosphor chromaticities on CIE diagram = displayable color range \\
\bottomrule
\end{tabular}
\end{center}
\begin{notesbox}
\textbf{Exam-day takeaway:} this whole cluster reduces to \textbf{five short, mechanical recipes} — none requires memorizing a proof, just the formula and the plug-in steps. The interpolation and matrix-$M$ questions are pure arithmetic once you know the setup; Phong and color-gamut are one-paragraph explain answers built directly from \S1.1 and \S1.6.
\end{notesbox}

\end{document}
